\subsection{局部化理论}
	\label{ss:locaff}
	\ssszh 我们想讨论局部几何朗兰兹的想法与技巧在仿射李代数表示论中的一个重要
	应用，而这一问题是纯代数的。先交代背景。

	\ssszh 回忆第 \ref{s:bbloc} 节讨论了
	约化李代数的局部化理论，并概述了如何用它确定
	单最高权模的特征标。

	对仿射李代数，Kashiwara--Tanisaki 在一系列工作中于非临界水平发展了类似的纲领
	\cite{kashiwara90}、\cite{kashiwaratanisakiposlevelii}、\cite{kashiwaratanisakinegativeintegrallevel}、
	\cite{kashiwaratanisakinegativerationallevel}、\cite{kashiwaratanisakipositiverational}、
	\cite{kashiwaratanisakialllevels}。扼要地说，在负水平，他们把仿射旗流形上的 D-模
	与某些 Kac--Moody 表示等同起来。在正水平，他们做了同样的事，但把仿射旗流形换成
	厚仿射旗流形，后者是无限型的概形，参数化 $\mathbb{P}^1$ 上在零点处有完整水平结构、在无穷远处有 Borel
	约化的 $G$-丛。作为推论，他们得到了非临界水平单最高权模的特征标。

	在他们工作之后，人们认识到，要推进剩下的情形，即临界水平的表示论，还需要新的想法与工具。下面我们
	想按 Beilinson--Drinfeld、Frenkel--Gaitsgory 与 Raskin--Yang 的工作，概述这里的局部化理论是什么样子。

	\ssszh 临界水平处的基本新困难在于：局部化定理的几何侧不再是某个空间上的 D-模，而是它的一种修正。
	令人愉快的是，只用我们迄今讨论过的局部几何朗兰兹内容，就能得到正确的陈述。下面来看怎么做到。

	首先，回忆对任一 $\cG$-oper $\chi$，设其底层的局部系统为 $\sigma$，局部几何朗兰兹预测：在映射
	\[
	  \sigma: \pt \rightarrow \LS
	\]
	之下对应于 $\Vect \simeq \QCoh(\pt)$ 的范畴环路群表示 $\sC_{\sigma}$，可以由重整化
	%
	\[
	  \fgk\mod \underset{\QCoh(\Op)} \otimes \Vect
	\]
	%
	得到，也就是取中心特征为 $\chi$ 的 Kac--Moody 表示。

   \begin{re} 既然我们已经讨论了范畴的归纳凝聚层，这一重整化可以理解为映射
   \[
   \Op \rightarrow \LS
   \]
   形式光滑的推论。特别地，这意味着由 $\QCoh(\Op)$ 生成的子 2-范畴只由拟凝聚范畴层构成。
   也就是说，没有空间容纳非平凡的奇异支撑。

   在朗兰兹对偶侧，这预测 $\fgc\mod$ 是缓增的，即典范映射
   %
   \[
        \Dmod(LG/ LN, \psi ) \underset {\Dmod(LN, \psi \backslash LG / LN,
        \psi)} \otimes \mathscr{W}_{\kappa_c}\mod \rightarrow \fgc\mod
   \]
   %
   是等价。这是关于临界水平 Kac--Moody 表示的一个值得注意的猜想。
   \end{re}

	\ssszh 把上述讨论（它对任意局部系统 $\sigma$ 都适用）特化到平凡联络。为此，可以平凡地把
	$\sigma$ 分解为
	%
	\[
	  \pt \rightarrow \pt/\cG \rightarrow \LS.
	\]
	%

	我们已经相当详细地讨论过 $\QCoh(\pt/\cG)$ 在局部几何朗兰兹下应当对应什么，即仿射
	Grassmann 簇上的 D-模。这给出 $\sC_{\sigma}$ 的另一种呈现，如下。由于 $\pt / \cG \rightarrow \LS$
	尽可能不光滑，$\Vect$ 落在由 $\QCoh(\pt / \cG)$ 生成的范畴之中。于是得到

	%
	{\footnotesize
	\begin{align*}
	     \Vect &\simeq  \QCoh( \pt / \cG) \\
	     &\quad \otimes_{\Hom_{\TwoIC(\LS)}(\QCoh(\pt / \cG), \QCoh(\pt / \cG))} \\
	     &\qquad \Hom_{\TwoIC(\LS)}( \QCoh(\pt / \cG), \QCoh(\pt)) \\
	     &\simeq
	     \QCoh( \pt / \cG) \otimes_{\ICn(\pt / \cG \underset \LS \times \pt /
	     \cG)} \ICn( \pt \underset \LS \times \pt ).
	\end{align*}}
	%
因此在自守侧，应当有
\begin{align}
   \sC_{\sigma} &\simeq \Dmod(\Gr_G)\underset {\ICn(\pt / \cG \underset \LS
   \times \pt /
   	\cG)} \otimes \ICn( \pt \underset \LS \times \pt ).
\intertext{而从 $\ICn( \pt / \cG \underset \LS \times \pt / \cG)$ 张量到
$\ICn(\pt
\underset  \LS \times \pt)$ 不过是忘掉 $\cG$ 作用的等变性，所以可以把它更直白地改写成 }
 \sC_{\sigma} & \simeq \Dmod(\Gr_G) \underset{\QCoh(\pt / \cG)} \otimes
 \QCoh(\pt).\label{e:heckeeig}
	\end{align}

显式地说，这就是仿射 Grassmann 簇上的 Hecke 特征层范畴。
也就是说，回忆 $\QCoh( \pt / \cG)$，即 $\on{Rep}(\cG)$，等同于阿贝尔 Satake 范畴
\[
\Dmod(L^+G \backslash LG /
L^+G)^\heartsuit.
\]
该范畴通过右卷积作用在仿射 Grassmann 簇的 D-模上，这一作用记为
%
\[
    - \overset {L^+G} \star -: \Dmod(LG/L^+G) \otimes \Dmod(L^+G \backslash LG
    / L^+G) \rightarrow \Dmod(LG / L^+G).
\]
%
所要的范畴由这一作用的特征对象构成，也就是说，由 $\Dmod(\Gr_G)$ 中的层 $\sF$ 构成，
并带有相容的同构
%
\begin{equation} \label{e:heckeeigen}
    \sF \overset {L^+G }\star \on{Sat}(V) \simeq \sF \otimes \on{Oblv}(V),
    \quad V \in \on{Rep}(\cG),
\end{equation}
%
其中 $\on{Oblv}(V)$ 记 $V$ 的底向量空间。关于这样的 $\sF$——与任一 Satake 层作卷积都返回
$\sF$ 的若干拷贝的直和——有一个基本观察：$\sF$ 必须有无限维的支撑，特别地，它不可能是
$\Dmod(\Gr_G)$ 的紧对象。

在比较 $\sC_\sigma$ 的两种描述之前，再作两条帮助定位的评注。

\begin{re}  把 Satake 层拖过张量积 \eqref{e:heckeeig}，也就是说，考虑
%
\[
     (\sF \overset {L^+G} \star V) \otimes k \simeq \sF \otimes (V \cdot k)
     \simeq
     \sF \otimes \on{Oblv}(V),
\]
%
其中 $\cdot$ 记 $\QCoh(\pt/\cG)$ 在 $\QCoh(\pt)$ 上的作用，就得到 Hecke 特征层更接地的表述
\eqref{e:heckeeigen}。
	\end{re}

	\begin{re} 初次接触这些想法的读者，可能觉得前面讨论的一些细节稍显复杂。因此值得强调，
	也可以通过类比来预测这一点。也就是说，在 $p$-adic 情形，要从
	%
	\[
	  \Fun( G(\Fqt)/G(\Fq[[t]]))
	\]
	%
	中提取 $G(\Fqt)$ 的一个不可约非分歧表示 $\pi_{\eta}$，人们取球 Hecke 代数右作用的特征向量，
	特征值为 $\eta$，即取相应的同态
	%
	\[   \eta: \Fun(G(\Fq[[t]]) \backslash G(\Fqt)/G(\Fq[[t]]))) \rightarrow k.\]
	%
	 在几何设定下，由 Tannaka 对偶，可能的特征值只有一个，即幺半函子
	 %
	 \[ \on{Oblv}: \on{Rep}(\cG) \rightarrow \Vect,\]
	 %
	  而我们再次是从万有表示出发、通过取特征向量来构造不可约非分歧表示。特别地，上面关于
	 Hecke 特征层支撑的观察，提升为关于 Hecke 特征向量支撑的类似断言。
	\end{re}

	\ssszh 比较非分歧情形下 $\sC_{\sigma}$ 的两种构造，就得到如下预测。也就是说，对平凡联络上的
	任一 oper 结构 $\chi$，应当有等价
	%
	\begin{equation} \label{e:locaff}
	      \Dmod(\Gr_G) \underset {\QCoh(\pt / \cG)} \otimes \QCoh(\pt) \simeq
	      \fgc\mod_\chi,
	\end{equation}
	%
	其中右端记中心特征为 $\chi$ 的 Kac--Moody 表示范畴的适当重整化。

	事实上，这是一条重要的定理，\footnote{更仔细地说，至少当 $\chi$ 对应形式穿孔圆盘上的 oper 时，
	它是定理，参见例 \ref{e:regop}。不过人们预期一般情形可以类似处理。}它的发展历时数十年，
	涉及多组作者的重要贡献。简要的历史如下。在对 Hitchin 系统的量子化的工作 \cite{bdh} 中，
	Beilinson--Drinfeld 证明了：在卷积作用
	%
	\[
	     \Dmod(L^+G \backslash LG / L^+G) \otimes \fgc\mod^{L^+G} \rightarrow
	     \fgc\mod^{L^+G}
	\]
	%
	之下，真空模的任一中心商都是阿贝尔 Satake 范畴作用的特征对象。这条定理是他们在整体几何朗兰兹中
	构造对应于整体 oper 的 Hecke 特征层时所用的主要局部输入。

	\ssszh Beilinson--Drinfeld 的工作对临界水平局部化理论的意涵，由 Frenkel--Gaitsgory 认识到，
	并在 \cite{frenkelgaitsgorylocgrassmannian04}、\cite{fg06loc} 与 \cite{fgloc09} 中大幅发展。首先，
	Frenkel--Gaitsgory 用 Beilinson--Drinfeld 建立的 Hecke 特征性质构造了一个函子
	%
	\begin{equation} \label{e:globgr}
	      \Dmod(\Gr_G)\underset {\QCoh(\pt / \cG)} \otimes \QCoh(\pt)
	      \rightarrow
	      \fgc\mod_\chi.
	\end{equation}
	%
	粗略地说，这相当于说明：$\Gr_G$ 上任一 D-模（以临界线丛扭之，从而得到临界水平的表示，并
	施加对应于 $\chi$ 的中心商）的整体截面自然地是 Hecke 特征对象。对 delta D-模，这正是
	Beilinson--Drinfeld 前述的定理，而一般情形可以化归到它。

	 除猜想 \eqref{e:globgr} 应是 $t$-正合的等价之外，Frenkel--Gaitsgory 还确立了许多基本性质。也就是说，
	 他们证明了它右 $t$-正合，且在下有界导出范畴上是全忠实的嵌入。\footnote{一般情形下的全忠实性，
	 特别是 Kac--Moody 表示范畴的正确重整化，后来由 Raskin \cite{RaskinLocCrit} 得到。}他们进一步证明了
	 它在 $I$-等变对象、即最高权模上是等价。

	\begin{re} 稍作加工，就可以从这些结果得到临界水平单最高权模特征标的 Kazhdan--Lusztig 猜想的类比。
	这一猜想由 Feigin--Frenkel 提出，见 \cite{arakawafiebig}，并在 \cite{DhillonYang} 中得到证明。
	\end{re}

\ssszh 在 Frenkel--Gaitsgory 的工作之后，其猜想剩下的部分是：\eqref{e:globgr} 是 $t$-正合的，且本质满。
后者可以有益地重新解释为环路群范畴作用下的一条生成性陈述，如下。\eqref{e:globgr} 的左端按定义由
其 $L^+G$-等变对象生成。因此，只需说明右端也如此。

一方面，这听起来是相当自然的陈述。确实，所施加的中心特征是非分歧性质的，对应形式非穿孔圆盘上的
oper，因此范畴本身应当是非分歧的，这颇为合理。然而，这一断言结果很难，很大程度上是因为人们对
驯分歧之外的 Kac--Moody 表示了解得很少，而且至今仍然如此。

\ssszh 十多年后，Raskin \cite{RaskinLocCrit} 取得了重要进展：他在第一个非平凡情形 $GL_2$
证明了该猜想，事实上对半单秩为一的任何群都证明了。为此，他证明了关于 $LGL_2$ 的范畴表示的一条
一般性断言，即它们由 Whittaker 不变量与 Iwahori 不变量生成；这与 $p$-adic 理论中一条更早、
同样基本的断言相平行。既然 \eqref{e:globgr} 在这类对象上是等价，完整的等价便随之而来。

猜想的一般情形随后由 Raskin--Yang \cite{RaskinYang} 解决。主要的新要素是 Yang \cite{yangmoyprasad}
把 Moy--Prasad 理论从 $p$-adic 情形移植过来；另见更早的 \cite{chenkamgarpour}。关键的是，除按深度
滤过环路群表示之外，Moy--Prasad 理论还为每个相继的商给出了一份紧生成元的显式清单，对应 $p$-adic
理论中的非精化极小 $K$-型。他们的证明随后通过一项非平凡的分析说明：从深度零以上的任一这样的生成元
出发若有非零映射，就与非分歧的中心特征不相容；于是得到所要的本质满性。

\begin{re} 事实上，前面讨论中我们略过了一个关键的技术要点：由于重整化，范畴 $\fgc\mod_\chi$ 并不
先验地带有 $LG$ 的范畴作用。这是 \cite{RaskinYang} 中克服的基本困难之一。
\end{re}

\begin{re} 我们想就 Moy--Prasad 理论在几何朗兰兹中的地位作一条评注，它与我们对局部化理论的讨论
本身无关。也就是说，由 Bremer--Sage \cite{sage1}、\cite{sage2}、\cite{sage3}，与 Galois 表示相对，
使用 $\cG$-联络有一个有趣的特点：对 $\cG$-联络也有极小 $K$-型理论。它与局部几何朗兰兹、特别是
与自守侧 Moy--Prasad 理论的相互作用，目前尚未被完全理解。	\end{re}

\ssszh 最后用两点评注结束关于临界水平局部化理论的讨论。第一，读者可以想见，上述内容可以从
非分歧情形推广到驯分歧情形，例如幂幺单值的情形。这由 Frenkel 与 Gaitsgory 在 \cite{fg06loc} 与
\cite{fglocaff} 中发展，做法大致是从仿射 Grassmann 簇过渡到仿射旗簇。

第二，颇为引人注目的是：尽管我们讨论的临界水平局部化定理已经很不平凡，它在某种意义上只是冰山
一角。也就是说，人们预期 $p$-adic 群更深层次的不可约表示的许多已知构造都应有几何版本，从而得到
等价，特别是带有野分歧中心特征的临界水平表示的局部化定理；相关讨论与猜想可参见 \cite{fg06loc} 与
\cite{kamgarpourschedler}。

特别地，综述作者预期：若一个 oper 的底层局部系统带有 $\cB$-约化，则相应的 Kac--Moody 模范畴
局部化为半无限旗流形上的扭 D-模范畴。对分歧特征标，这样的局部化不再能像 \cite{fglocaff} 那样在
仿射旗簇上实现，而要在先由 Roche \cite{roche} 在 $p$-adic 情形引入的某个更小的紧开子群的商上实现。
